
\chapter{Notation and Preliminaries}\label{ch:notation}

\section{Notation}
For a positive integeger $n$ define $[n] = \Brace{1,\ldots,n}$.
We use $\lambda \in \N$ exclusively to denote the security parameter.
We use $\poly$ to denote the set of functions that are polynomial in $\lambda$
in the sense of \Cref{def:polynomial} and $\negl$ to denote the set of functions that are 
negligible in $\lambda$ in the sense of \Cref{def:negligible}.
We call a function $f$ \textit{overwhelming} if $1 - f(\lambda) \in \negl$.
We use $\gets$ for simple assignment.
For a set $S$ we write $s \sample S$ to indicate that the value $s$ is sampled uniformly at random from 
set $S$ and for a probabilistic algorithm $\A$ we write $y \sample \A(x)$ to denote assigning to 
$y$ the output of $\A$ on input $x$. We use ppt. to abbreviate probabilistic polynomial time in the context of 
Turing machines. Finally, we use bold letters, e.g. $\x$, to indicate that $\x$ consists of multiple values such as tuples or vectors\todo{maybe change}.

\begin{definition}[Polynomial Functions]\label{def:polynomial}
	We say that a function $p: \N \to \R$ is polynomial if the folowing holds

	\[
		\exists d \in \N, \exists n_0 \in \N, \forall n > n_0, p(n) \le n^d
	\]
\end{definition}

\begin{definition}[Negligible Functions]\label{def:negligible}
	We say that a function $\eps: \N \to \R$ is negligible if the following holds

	\[
		\forall p \in \poly, \exists n_0 \in \N, \forall n > n_0, |\eps(n)| < \frac{1}{p(n)}
	\]
\end{definition}

\newpage

\section{Verifiable Functions}
A verifiable function (VF) $\prim{VF}$ is a cryptographic primitive that implements algorithms 
$\alg{Gen}{VF}, \alg{Eval}{VF}$ and $\alg{Vfy}{VF}$. A user can generate key-pairs $(\sk, \pk)$
using $\alg{Gen}{VF}$ consisting of a secret key $\sk$ and a public key $\pk$, use the secret key $\sk$
to make evaluate arbitrary points $x \in \bin^\lambda$ using $\alg{Eval}{VRF}$ and use the public key $\pk$ to 
verify a triple $(x,y,\pi)$ using $\alg{Vfy}{VRF}$.
We give a formal definition in \Cref{def:vf}, the definition takes inspiration of the definition of a 
VRF in \cite{}. To reduce redundancy we define additional security properties 
such as pseudorandomness and unpredictability seperately.

\begin{definition}[Verifiable Function]\label{def:vf}
	A verifiable function (VF) $\prim{VF}$ is a triple of efficient (ppt.) algorithms
	$\Paren{\alg{Gen}{VF}, \alg{Eval}{VF}, \alg{Vfy}{VF}}$ with the syntax\todo{this is less elegant now}

	\begin{itemize}
		\item $\alg{Gen}{VF}$: The key-generation algorithm takes as input the unary encoding of the security parameter
			and produces a pair $(\sk, \pk) \in \SK \times \PK$ containing a secret and a public key respectivley.
		\item $\alg{Eval}{VF}$: The evaluation algorithm takes as input a tuple $(\sk, x)$ containing a secret key $\sk \in \SK$ and an evaluation point $x \in \bin^\lambda$
			and produces an tuple $(y, \pi)$ containing an image $y \in \Y$ and a proof $\pi \in \P$. 
		\item $\alg{Vfy}{VF}$: The verification algorithm takes as input a tuple $(\pk, x, y, \pi)$ containing a public key $\pk \in \PK$, an evaluation
			point $x \in \bin^\lambda$, an image $y \in \Y$ and a proof $\pi \in \P$ and returns a bit $b$ denoting acceptance or rejection.
	\end{itemize}
\end{definition}

\todo{can evaluation or verification be randomized, probably vfy can be probabilistic}
We use the notation $\alg{Eval}{VF}^y$ and $\alg{Eval}{VF}^\pi$ to denote the 
algorithm $\alg{Eval}{VF}$ to projected on the respective output.
Note that a VF is a purely syntactical concept. For a VF we define various properties to arrive to arrive 
at specific cryptographic primitives, e.g. VRFs. First of all we define $\err$-correctness, 
a property that guarantes that over random key-pairs and evaluation points, the 
verification algorithm accepts a correctly computed image-proof pair with probability at least $\err$.


\begin{definition}[$\err$-Correctness, \cite{}]\label{def:correctness}
	A verifiable function $\prim{VF} = (\alg{Gen}{VF}, \alg{Eval}{VF}, \alg{Vfy}{VF})$ is 
	$\err$-correct, if for all security parameters $\lambda \in \N$ we have that

	\[
		\Prob*{
			\alg{Vfy}{VF}(\pk, x, y, \pi) = 1 : 
			\begin{aligned}
				(\sk, \pk) &\sample \alg{Gen}{VF}(1^\lambda)\\ 
				x &\sample \bin^\lambda\\
				(y, \pi) &\gets \alg{Eval}{VF}(\sk, x)
			\end{aligned}
		} \ge \err(\lambda)
	\]

	We define a VF $\prim{VF}$ to be perfectly correct, or just correct, if it is $1$-correct and statistically correct if 
	$\err$ is an overwhelming function.
\end{definition}

Intuitively a correct VF states that honestly generated image-proof pairs are accepted by the 
verification algorithm. Next we define unique provability. A VF is uniquely provable if 
not more than one image can be proven, even for maliciously crafted public keys.

\begin{definition}[Unique Provability, \cite{}]\label{def:unique_provability}
	A verifiable function $\prim{VF}$, with algorithms $(\alg{Gen}{VF}, \alg{Eval}{VF}, \alg{Vfy}{VF})$,
	is defined to be uniquely provable if for all evaluation points $x \in \bin^\lambda$, all public-keys $\pk \in \PK$
	and any two images-proof pairs $(y_0,\pi_0), (y_1,\pi_1) \in \Y \times \PK$ we have:

	\[
		\alg{Vfy}{VF}(\pk, x, y_0, \pi_0) = 1 \land \alg{Vfy}{VF}(\pk, x, y_1, \pi_1) = 1 \implies y_0 = y_1
	\]
\end{definition}

\begin{comment}
\begin{definition}[$\err$-Unique Provability]\label{def:unique_provability}
	A verifiable function $\prim{VF}$, with algorithms $(\alg{Gen}{VF}, \alg{Eval}{VF}, \alg{Vfy}{VF})$,
	defined to be $\err$-uniquely provable if over 
	the randomness of the key generation algorithm and for all evaluation points $x \in \bin^\lambda$, the probability of there existing $(y_0, \pi_0), (y_1, \pi_1) \in \Y \times \P$
	with $y_0 \neq y_1$ such that both verify is at most $\err$, i.e. if
	\[
		\Prob*{
			\begin{aligned}
				&\exists (y_0, \pi_0), (y_1, \pi_1) \in \Y \times \P, y_0 \neq y_1 \\
				&\land \alg{Vfy}{VF}(\pk, x, y_0, \pi_0) = 1\\ 
				&\land \alg{Vfy}{VF}(\pk, x, y_1, \pi_1) = 1
			\end{aligned}
			:
			\begin{aligned}
				x  &\in \bin^\lambda\\ 
				(\sk, \pk) &\sample \alg{Gen}{VRF}(1^\lambda)
			\end{aligned} 
		} \le \err
	\]
	key-pairs sampled according to the key generation algokorithm 
	the  for all evaluation points $x \in \bin^\lambda$, a public-keys $\pk \in \PK$
	(even maliciously crafted ones) and any two images-proof pairs $(y_0,\pi_0), (y_1,\pi_1) \in \Y \times \PK$:

	\[
		\alg{Vfy}{VF}(\pk, x, y_0, \pi_0) = 1 \land \alg{Vfy}{VF}(\pk, x, y_0, \pi_0) = 1 \implies y_0 = y_1
	\]
\end{definition}
\end{comment}

\todo{maybe comment why no statistical notion}

Correctness and unique provability are information theoretic. We define pseudorandomness of a VF
which is a computational property, stating that no efficient adversary $\A$ can distinguish images $y$ 
corresponding to evaluation of $x$, chosen by $\A$, from a uniformly random output $y^* \sample \Y$.
This needs to hold even if the adversary gets to see the evaluation of a polynoial number of points $x \neq x^*$.

\begin{definition}[$\errtime$-Pseudorandomness, \cite{}]\label{def:pseudorandomness}
	A VF $(\alg{Gen}{VF}, \alg{Eval}{VF}, \alg{Vfy}{VF})$ is defined to be $\errtime$-pseudorandom if for all security parameters $\lambda \in \N$, 
	for all stateful adversaries $\A$, running in time at most $\runtime(\lambda)$, the advantage of $\A$ in the prf security experiment outlined in \Cref{fig:prf_wprf_up_experiment} 
	is at most $\err(\lambda)$, i.e. if 

	\[
		\Adv{VF}{prf}{\lambda, \A} := 2 \cdot \abs*{\Prob*{\Exp{VF}{prf}(\lambda, \A) = 1} - \frac{1}{2}} \le \err(\lambda)
	\]

	We say that a VF is pseudorandom, if it is $\errtime$-pseudorandom with $\err$ being a negligible function and $\runtime$ being a polynomial, i.e. $\err \in \negl, \runtime \in \poly$.
\end{definition}

Next we define two weaker notions of pseudorandomness. While (strong) pseudorandomness required the evaluation of adverserially chosen points 
to be indistinguishable from uniformly random elements $y^* \in \Y$ we can relax this requirement by requiring computational indistinguishability
for uniformly random challenges $x^* \sample \bin^\lambda$. We call this notion weak pseudorandomness.

\begin{definition}[$\errtime$-Weak Pseudorandomness, \cite{}]\label{def:weak_pseudorandomness}
	A VF $(\alg{Gen}{VF}, \alg{Eval}{VF}, \alg{Vfy}{VF})$ is defined to be $\errtime$-weakly pseudorandom if for all security parameters $\lambda \in \N$, 
	for all stateful adversaries $\A$, running in time at most $\runtime(\lambda)$, the advantage of $\A$ in the wprf security experiment outlined in \Cref{fig:prf_wprf_up_experiment} 
	is at most $\err(\lambda)$, i.e. if 

	\[
		\Adv{VF}{wprf}{\lambda, \A} := 2 \cdot \abs*{\Prob*{\Exp{VF}{wprf}(\lambda, \A) = 1} - \frac{1}{2}} \le \err(\lambda)
	\]

	We say that a VF is weakly pseudorandom, if it is $\errtime$-weakly pseudorandom with $\err$ being a negligible function and $\runtime$ being a polynomial, i.e. $\err \in \negl, \runtime \in \poly$.
\end{definition}

Alternatively a VF can be unpredictable. A VF is unpredictable if no adversary 
can predict the image $y$ of the evaluation of a challenge $x$ chosen by the adversary.
This property needs to hold even if the adversary sees the evaluation of a polynomial number of points $x \neq x^*$.
Note that if the adversary can predict the output of evaluation on $x^*$ he can also distinguish the output from a random output $y^* \sample \bin$
with non-negligible probability making unpredictability is a weaker property than pseudorandomness.

\begin{definition}[$(\eps, t)$-Unpredictability, \cite{}]\label{def:unpredictability}
	A VF $(\alg{Gen}{VF}, \alg{Eval}{VF}, \alg{Vfy}{VF})$ is defined to be $\errtime$-unpredictable if for all security parameters $\lambda \in \N$, 
	for all stateful adversaries $\A$, running in time at most $\runtime(\lambda)$, the advantage of $\A$ in the up security experiment outlined in \Cref{fig:prf_wprf_up_experiment} 
	is at most $\err(\lambda)$, i.e. if 

	\[
		\Adv{VF}{up}{\lambda, \A} := 2 \cdot \abs*{\Prob*{\Exp{VF}{up}(\lambda, \A) = 1} - \frac{1}{2}} \le \err(\lambda)
	\]

	We say that a VF is unpredictable, if it is $\errtime$-unpredictable with $\err$ being a negligible function and $\runtime$ being a polynomial, i.e. $\err \in \negl, \runtime \in \poly$.
\end{definition}

\begin{figure}[htbp]
	\begin{pcvstack}[center, space=1em]
	\begin{pchstack}[space=1em]
		\procedure[linenumbering]%
		{Experiment $\Exp{VF}{prf}(\lambda, \A)$}
		{
			b \sample \bin\\ 
			(\sk, \pk) \sample \alg{Gen}{VF}(1^\lambda)\\ 
			S \gets \varnothing\\ 
			\highlight{(x^*, \sigma) \sample \mathcal A^{\O_{\text{eval}}}(1^\lambda, \pk)}\\ 
			(y_0, \pi_0) \gets \alg{Eval}{VF}(\sk, x^*)\\ 
			y_1 \sample \Y\\ 
			b' \sample \A(1^\lambda, x^*, \sigma, y_b) \\ 
			\pcreturn b = b' \land x^* \notin S
		}
		\procedure[linenumbering]%
		{Experiment $\Exp{VF}{wprf}(\lambda, \A)$}
		{
			b \sample \bin\\ 
			(\sk, \pk) \sample \alg{Gen}{VF}(1^\lambda)\\ 
			S \gets \varnothing\\ 
			\sigma \sample \A^{\O_{\text{eval}}}(1^\lambda, \pk)\\ 
			\highlight{x^* \sample \bin^\lambda}\\ 
			(y_0, \pi_0) \gets \alg{Eval}{VF}(\sk, x^*)\\ 
			y_1 \sample \Y\\ 
			b' \sample \A^{\O_{\text{eval}}}(1^\lambda, x^*, \sigma, y_b) \\ 
			\pcreturn b = b' \land x^* \notin S
		}
	\end{pchstack}
	\begin{pchstack}[space=1em]
		\procedure[linenumbering]%
		{Experiment $\Exp{VF}{up}(\lambda, \A)$}
		{
			(\sk, \pk) \sample \alg{Gen}{VF}(1^\lambda)\\ 
			S \gets \varnothing\\ 
			(x^*, y^*) \sample \A^{\O_{\text{eval}}}(1^\lambda, \pk)\\ 
			(y, \pi) \gets \alg{Eval}{VF}(\sk, x^*)\\ 
			\pcreturn y^* = y \land x^* \notin S
		}
		\procedure[linenumbering]%
		{Oracle $\O_{\text{eval}}(x)$}
		{
			S \gets S \cup \{x\}\\ 
			(y,\pi) \sample \alg{Eval}{VF}(\sk, x)\\ 
			\pcreturn y
		}
	\end{pchstack}
	\end{pcvstack}
	\caption{The pseudorandomness security experiment $\Exp{VF}{prf}(\lambda, \A)$ used in \Cref{def:pseudorandomness}, 
	the weak pseudorandomness security experiment $\smash{\Exp{VF}{wprf}(\lambda, \A)}$ used in \Cref{def:weak_pseudorandomness} and
	the unpredictability security experiment $\Exp{VF}{up}(\lambda, \A)$ used in \Cref{def:unpredictability}.
	The main changes between the pseudorandomness security experiment and the weak pseudorandomness 
	security experiment are highlighted in grey.}
	\label{fig:prf_wprf_up_experiment}
\end{figure}

We now define the notions of a verifiable random function (VRF), verifiable unpredictable function (VUF)
and weak verifiable random function (WVRF), which are correct, uniquely provable and 
pseudorandom, unpredictable and weakly pseudorandom VFs respectivly.
We formally define the primitives.

\begin{definition}[Verifiable Random Functions]\label{def:vrf}
	A verifiable random function $\prim{VRF} = (\alg{Gen}{VRF}, \alg{Eval}{VRF}, \alg{Vfy}{VRF})$ is a VF that satisfies 
	perfect correctness in the sense of \Cref{def:correctness}, unique provability in the sense of \Cref{def:unique_provability} and pseudorandomness
	in the sense of \Cref{def:pseudorandomness}.
\end{definition}

\begin{definition}[Weak Verifiable Random Functions]\label{def:wvrf}
	A weak verifiable unpredictable function $\prim{WVRF} = (\alg{Gen}{WVRF}, \alg{Eval}{WVRF}, \alg{Vfy}{WVRF})$ is a VF that satisfies 
	perfect correctness in the sense of \Cref{def:correctness}, unique provability in the sense of \Cref{def:unique_provability} and weak pseudorandomness
	in the sense of \Cref{def:weak_pseudorandomness}.
\end{definition}

\begin{definition}[Verifiable Unpredictable Functions]\label{def:vuf}
	A verifiable unpredictable function $\prim{VUF} = (\alg{Gen}{VUF}, \alg{Eval}{VUF}, \alg{Vfy}{VUF})$ is a VF that satisfies 
	perfect correctness in the sense of \Cref{def:correctness}, unique provability in the sense of \Cref{def:unique_provability} and unpredictability
	in the sense of \Cref{def:unpredictability}.
\end{definition}

\begin{comment}
Unbiasability, following the definition of \cite{EC:GiuSte24}, is a property for VRFs that 
prevents adversaries from skewing the output distribution of the evaluation function on malleable keys. 
To formally definine unbiasability we require the notion of monotone predicates.

\begin{definition}[Monotone Predicates \cite{}]

	A probabilistic predicate $\rho$ on the domain $(\Y \cup \{\bot\})^*$ is monotone, iff 
	for all $y \in \Y^*$, all index sets $I \subseteq [|y|]$ and every bitstring of randomness $r \in \bin^\lambda$ we have

	\[
		\Prob*{\rho(\tilde y; r) = 1} \le \Prob*{\rho(y; r) = 1}
	\]

	where $\tilde y_i = y_i$ if $i \notin I$ and $\tilde y_i = \bot$ otherwise.
\end{definition}

\begin{definition}[Unbiasability \cite{}]
	A VRF $\prim{VRF}$ is (fully) $\errtimelen$-unbiasable if for every stateful adversary $\A$ running in time 
	at most $\runtime$ its advantage in the unbiasability security experiment in \Cref{fig:fub_experiment} is bounded by 

	\[
		\Adv{VRF}{fub}{\lambda, \A} := 2\cdot \Prob*{\Exp{VRF}{fub}(\lambda, \mathfrak{a}, \mathfrak{b}, \A) = 1} - 1 \le \err(\lambda)
	\]

	We say that a VRF is unbiasable if it is $\errtimelen$-unbiasable and $\err$ is a negligible function in the 
	security parameter $\lambda$, i.e. $\err \in \negl$.
\end{definition}

\begin{figure}[hptb]
	\begin{pchstack}[center, space=1em]
		\procedure[linenumbering]%
		{Experiment $\Exp{VF}{fub}(\lambda, \A)$}
		{
			b \sample \bin\\ 
			(\rho, \pk_1, \ldots, \pk_{\mathfrak{a}(\lambda)}) \sample \A(1^\lambda)\\ 
			\textbf{require }\rho \text{ is monotone}\\ 
			x_1, \ldots, x_{\mathfrak{b}(\lambda)} \sample \bin^\lambda\\ 
			(y_{1,1}^0, \pi_{1,1}^0), \ldots, (y_{\mathfrak{a}(\lambda), \mathfrak{b}(\lambda)}^0, \pi_{\mathfrak{a}(\lambda), \mathfrak{b}(\lambda)}^0) \sample \A(x_1, \ldots, x_{\mathfrak{b}(\lambda)})\\ 
			\pcfor i \in [\mathfrak{a}(\lambda)], j \in [\mathfrak{b}(\lambda)] \pcdo\\ 
			\t \beta_{i,j} \gets \alg{Vfy}{VRF}(\pk_i, x_j, y_{i,j}^0, \pi_{i,j}^0)\\ 
			\t \pcif \beta_{i,j} = 1 \pcthen\\ 
			\t \t z_{i,j} \gets y_{i,j}^b\\ 
			\t \pcelse z_{i,j} \gets \bot\\ 
			r \gets \bin^\lambda\\ 
			\pcreturn \rho(z_{1,1}, \ldots, z_{\mathfrak{a}(\lambda), \mathfrak{b}(\lambda)}, r) \neq b
		}
	\end{pchstack}
	\caption{Full unbiasability security experiment.}
	\label{fig:fub}
\end{figure}

Key-homomorphism of VF requires that the secret-key space and the product space of the image and proof space admit group structures $(\SK, \odot)$ and $(\Y \times \P, \otimes)$ respectively.
Intuitively a VF is key-homomorphic if for every evaluation point $x \in \bin^\lambda$ the evaluation algorithm on $x$ implements a grouphomomorphism in the secret key.

\begin{definition}[Key-Homomorphism]
	A VF $\prim{VF}$ is key-homomorphic if for every evaluation point $x \in \bin^\lambda$, secret key $\sk_1, \sk_2 \in \SK$ we have 

	\[
		\alg{Eval}{VF}(\sk_1 \odot \sk_2, x) = \alg{Eval}{VF}(\sk_1, x) \otimes \alg{Eval}{VF}(\sk_2, x) 
	\]
\end{definition}

We also refer to key-homomorphism as secret-key aggregability. To define public-key aggregability we first define
the cryptoraphic primitive of an aggregable verifiable function (AggVF) that additionally implements algorithms to combine 
image-proof pairs into a single image-proof pair aggregate.

\begin{comment}
Aggregatability \cite{} is a property on VRFs that wllows the aggregation of either secret-keys for 
secret-key aggregatability or image-proof pairs for public-key aggregatability. Secret-key aggregatability, also called key-homomorphism, 
requires the evaluation function to consitute a group homomorphism from a group structure on the secret key space $\SK$ 
to a group structure on the output codomain $\Y$. We follow \cite{} and formally define key-homomorphism.

\begin{definition}[Key-Homomorphism, \cite{}]
	A VRF $\prim{VRF}$ is key homomorphic if $(\Y \times \P, \otimes), (\SK, \odot)$ are groups and for all secret keys $(\sk_0, \sk_1) \in \SK$
	and all evaluation points $x \in \bin^\lambda$ we have 

	\[
		\alg{Eval}{VRF}(\sk_0, x) \otimes \alg{Eval}{VRF}(\sk_1, x) = \alg{Eval}{VRF}(\sk_0 \odot \sk_1, x)
	\]
\end{definition}

For public-key aggregatability we first define the syntax for a aggregatable VF. That is a VF with two additional 
algorithms, one for combining image-proof pairs into a single image-proof pair and one for verifying whether 
an aggregate image-proof pair was computed correctly.

\begin{definition}[Aggregatable Verifiable Functions]
	An aggregatable VF $\prim{AggVF}$ is a VF with an additional tuple of deterministic polynomial-time algorithms 
	with the following syntactic specification.

	\begin{itemize}
		\item $\alg{Agg}{AggVF}(x, \{\pk_i, y_i, \pi_i\}_{i \in [n]})$: The aggregation algorithm takes as input an evaluation point $x \in \bin^\lambda$,
			and $n$ public keys $\pk_1, \ldots, \pk_n \in \PK$ and $n$ image-proof pairs $(y_1, \pi_1), \ldots, (y_n, \pi_n)$, and 
			outputs an aggregate key $\widetilde \pk$ as well as an aggregate image-proof pair $(\widetilde y, \widetilde \pi)$.
		\item $\alg{AggVfy}{AggVF}(\{\pk_i\}_{i \in [n]}, x, \widetilde y, \widetilde y)$: The aggregate verification algorithm 
			takes as input $n$ public keys $\pk_1, \ldots, \pk_n \in \PK$, an evaluation point $x \in \bin^\lambda$, and an aggregate image-proof pair 
			$(\widetilde y, \widetilde \pi)$ and outputs a bit $b \in \bin$ indicating acceptance or rejection.
	\end{itemize}
\end{definition}

We define several properties for aggregatable verifiable functions (AggVF), following \cite{}.

\begin{definition}[$\err$-Aggregate Correctness,]
	An AggVF $\prim{AggVF}$ is $\err$-aggregate correct, if for all security parameters $\lambda \in \N$, key-pairs $(\sk_i, \pk_i) \sample \alg{Gen}{AggVF}(1^\lambda)$, 
	and evaluation points $x \in \bin^\lambda$ we have 
	
	\[
		\Prob*{
			\begin{aligned}
				(\widetilde{y}, \widetilde \pi) \sample \alg{Agg}{AggVF}(x, \{\pk_i, \alg{Eval}{AggVF}(\sk_i, x)\}_{i \in [n]})\\ 
				\alg{AggVfy}{AggVF}(\{\pk_i\}_{i \in [n]}, x, \widetilde y, \widetilde \pi) = 1
			\end{aligned}
		} \ge \err(\lambda)
	\]
\end{definition}
	
\begin{definition}[Aggregate Pseudorandomness]
	
\end{definition}

\begin{definition}[Aggregate Binding]
	
\end{definition}

\begin{definition}[Aggregate Uniqueness]
	
\end{definition}
\end{comment}

\newpage

\section{Robust Combiners}
For two natural numbers $n,k$ we define the robust $(n;k)$ black-box combiner, or robust black-box combiner for short.
A robust $(n;k)$ black-box combiner for a primitive $\prim{P}$ is a construction $\BBComb{P}$, to be precise a ppt. oracle Turing machine,
that is given query access to $n$ oracles, the candidate schemes $\O_1, \ldots, \O_n$. If at least $k$ candidate schemes implement 
the primitive $\prim{P}$, then the combiner $\BBComb{P}[\O_1, \ldots, \O_n]$ securely implements $\prim{P}$.
The idea behind a robust black-box combiner is that the construction extracts the security properties from the subset of 
correct candidates, even if some candidate schemes are not secure instances. However now note the following problem in the 
current definition:

If there exists a secure instance of $\prim{P}$ there always exists robust $(n;0)$ black-box combiner that simply 
ignores its candidates and implements $\prim{P}$. This robust black-box combiner meets all the requirements of the definition.
Showing that there exists no specific robust black-box combiner therefore requires to show that no secure instance 
of the primitive $\prim{P}$ exists, which is often a difficult complexity theoretic statement.

To circumvent this problem we want to explicitly force the combiner to derive its security from the 
subset of secure candidate schemes. To do this, we require the security properties of the combiner to hold 
relative to a PSPACE complete oracle. Firstly we want to note that this means that the adversary against the 
combiner can efficiently check whether an element is a member of a PSPACE complete language, enabling the adversary 
to generally break computational hardness assumptions but limiting the adversary still to a polynomial number of 
oracle queries. Secondly we want to note that this notion of a combiner still doesnt solve the issue 
if the security of the cryptographic primitive are \textit{all} information theoretic.

\begin{comment}
Given a cryptographic primitive $\prim{P}$ we define a robust $(n;k)$ $\prim{P}$-combiner

\begin{definition}[Robust $(n;k)$ $\prim{P}$-combiner]
	Given a cryptographic primitive $\prim{P}$ describing $p$ algorithms $({\alg{alg}{P}}^1, \ldots, {\alg{alg}{P}}^2)$,
	a robust $(n;k)$ $\prim{P}$-combiner is a $p$-tuple of algorithms that take as input the description of $n$ Turing machines $\O_1, \ldots, \O_n$, 
	the candidate schemes
	
	\[
		\Comb{P}[\O_1, \ldots, \O_n] = \Paren*{\Comb{P}[\O_1, \ldots, \O_n].{\alg{alg}{P}}^1, \ldots, \Comb{P}[\O_1, \ldots, \O_n].{\alg{alg}{P}}^n}
	\]

	that run in time polynomial in $n$, and the size of their input, 
	such that the combiner securely implements the primitive $\prim{P}$ if at least $k$ candidates schemes securely 
	implement the primitive $\prim{P}$.
\end{definition}
\end{comment}

% We write robust combiner for robust $(n;k)$ $\prim{P}$-combiner whenever $n,k$ and $\prim{P}$ are clear from the context. 
% A robust combiner gets the candidate schemes as a description of their turing machine and can therefore be 
% adaptive in the code of its candidates. Note that this definition allows the robust combiner to ignores its candidate schemes 
% and implement a secure instance of the underlying primitive indepdently of its candidates. One can therefore always construct 
% a robust combiner for primitives that admit secure instances, by contraposition this also means that any result 
% that rules out the existence of a robust combiner for a specific primitive $\prim{P}$ also rules out the existence 
% of the primitive itself. Intuitively we would like for the combiner to derive its security from the security of its 
% secure candidate schemes, therefore we define robust black-box combiner. A robust black-box combiner is a robust combiner 
% that, instead of getting access to the description of the Turing machines of the candidate schemes, only gets oracle access 
% to the underlying candidate schemes. We define the notion of a robust black-box combiner closely following the definition 
% by \cite{}. Note that we explicitly require a robust black-box combiner to derive its security from its candidates,
% for that reason we require the security to hold relative to a PSPACE oracle $\PSPACE$.

\todo{Very unhappy with this definition}

TODO state contrapositive of the definition below here.

\begin{definition}[Robust $(n;k)$ Black-Box $\prim{P}$-Combiner]\label{def:robust_blackbox_combiner}
	A robust $(n;k)$ black-box combiner for primitive $\prim{P}$ with $k$ algorithms is a tuple
	of ppt. oracle Turing machines $\BBComb{P} = (\alg{Alg}{P}^1, \ldots, \alg{Alg}{P}^{k})$
	such that if there exists an ppt. adversary $\A^{\O_1, \ldots, \O_n, \PSPACE}$ that against 
	$\BBComb{P}[\O_1, \ldots, \O_n]$ then there exist $n-k+1$ ppt. adversaries $\smash{\B_1^{\O_{\ell(1)}}, \ldots, \B_{n-k+1}^{\O_{\ell(n-k+1)}}}$
	and an injective function $\ell: [n-k+1] \to [n]$ such that for any $i \in [n-k+1]$ the adversary $\smash{\B_i^{\O_{\ell(i)}}}$ breaks the security of the candidate $\O_{\ell(i)}$.
\end{definition}

\paragraph{Remark.}
If $\prim{P}$ defines composite security properties, as is the case for VRFs for instance, 
we say that an adversary breaks the security of $\prim{P}$ if it breaks one of the required security properties.
For computational security properties such as pseudorandomness this typically means that the adversary 
wins the corresponding security experiment with non-negligible advantage. For information theoretic 
or statistical properties such as unique provability or correctness this means that the property doesn't hold
(we can think of this as the adversary $\A$ being arbitrary, breaking the security).